\chapter{Introducción}
\label{chap:intro}

\ph{Este capítulo presenta el trabajo. Da al lector el contexto mínimo para entender
qué se aborda, por qué es relevante, qué se pretende conseguir y qué queda fuera.
No se entra todavía en detalles técnicos: eso pertenece al estado del arte y al diseño.
Fija aquí el registro que mantendrás en toda la memoria: redacción impersonal
(``se ha desarrollado'') o primera persona del plural (``hemos desarrollado''),
pero consistente de principio a fin, y prosa académica sin coloquialismos
(consulta la Guía de estilo incluida al inicio del documento en modo borrador).}

\section{Contexto}
\label{sec:contexto}

\ph{Describe el dominio en el que se enmarca el TFG y por qué es relevante.
Aporta el trasfondo necesario para que un lector ajeno al tema entienda
el resto del documento. Un lector tras leer esta sección debería poder responder a:
``¿de qué va este trabajo y en qué entorno se sitúa?''.}

\ph{Sugerencia de contenido:
\begin{itemize}
  \item \ph{Sector o disciplina (industria, investigación, educación\dots).}
  \item \ph{Tendencias o cambios recientes que motivan el problema.}
  \item \ph{Conceptos básicos imprescindibles para entender el resto.}
\end{itemize}}

El desarrollo de software actual debe dar respuesta a sistemas cada vez más configurables, capaces de adaptarse a distintas necesidades, plataformas y contextos de uso. Dentro de este escenario, las Líneas de Productos Software (Software Product Lines, SPL) \cite{spl} ofrecen un enfoque para gestionar familias de productos relacionados a partir de los elementos que comparten y de aquellos que pueden variar entre unas configuraciones y otras. Esto permite tratar la variabilidad de forma explícita y facilita tareas como la configuración de productos, la reutilización de componentes o el análisis automático de las distintas combinaciones posibles.

Una de las formas más extendidas de representar esta variabilidad son los modelos de características (Feature Models) \cite{featureModels}. Estos modelos describen las características que puede incluir una familia de productos y las relaciones que existen entre ellas, permitiendo establecer qué elementos son obligatorios u opcionales, cuáles pueden combinarse y qué restricciones deben respetarse para obtener una configuración válida.

A medida que los modelos han ganado en complejidad, también ha surgido la necesidad de disponer de formatos que permitan representarlos de manera común. Universal Variability Language (UVL) \cite{uvl} nace con este propósito y proporciona un lenguaje para describir modelos de variabilidad con diferentes grados de expresividad. Además de las relaciones habituales entre características, UVL permite incorporar construcciones más avanzadas, como atributos y restricciones, ampliando así los escenarios que pueden representarse y analizarse.

En torno a estos modelos se han desarrollado distintas herramientas destinadas a facilitar su procesamiento y análisis. Entre ellas se encuentra Flamapy \cite{flamapy}, un framework desarrollado en Python que proporciona soporte para el análisis automático de modelos de características. Su arquitectura extensible permite incorporar nuevos formatos, operaciones y técnicas de análisis sin necesidad de modificar el núcleo de la herramienta, lo que facilita su evolución y su uso en diferentes líneas de investigación.

El presente Trabajo Fin de Grado se desarrolla dentro de este ecosistema y se enmarca en la actividad investigadora de DiversoLab y en TASOVA Plus (Red en nuevas Tendencias en Arquitectura Software y Variabilidad) \cite{tasovaplus}. Este entorno promueve el desarrollo de herramientas relacionadas con la ingeniería de líneas de productos software, la arquitectura software y el análisis de variabilidad, siguiendo además principios asociados al software libre y a la ciencia abierta (Open Science).

Otro de los componentes de este ecosistema es UVLHub \cite{uvlhub}, una plataforma orientada a la publicación, gestión y exploración de modelos expresados en UVL. La plataforma ofrece un espacio en el que investigadores y desarrolladores pueden compartir modelos y reutilizarlos en distintos experimentos o herramientas. Flamapy, UVL y UVLHub conforman así buena parte del entorno tecnológico sobre el que se desarrolla este trabajo.

\section{Problema}
\label{sec:problema}

\ph{Plantea el problema concreto que el TFG resuelve. Conviene formularlo en términos
de limitaciones reales del estado actual: qué no se puede hacer hoy o qué se hace mal,
y qué consecuencias tiene.}

\ph{Estructura recomendada (1--3 problemas claramente delimitados):}

\ph{\begin{enumerate}
  \item \textbf{\ph{Limitación 1:}} \ph{descripción del problema y de su impacto.}
  \item \textbf{\ph{Limitación 2:}} \ph{descripción del problema y de su impacto.}
  \item \textbf{\ph{Limitación 3:}} \ph{descripción del problema y de su impacto.}
\end{enumerate}}

La evolución de los lenguajes de modelado y de los motores de razonamiento ha permitido trabajar con formas de variabilidad cada vez más complejas. Sin embargo, disponer de herramientas capaces de procesar esa complejidad no resulta suficiente por sí solo: para poder evaluarlas también es necesario contar con modelos que hagan uso de esas capacidades.

Buena parte de los modelos de características disponibles se concentra en construcciones relativamente básicas. Estos modelos siguen siendo útiles para numerosos análisis, pero resultan menos adecuados cuando se pretende experimentar con elementos más avanzados del lenguaje o estudiar cómo se comportan los motores de razonamiento ante distintos grados de complejidad. Se produce así una diferencia entre la expresividad que pueden ofrecer los lenguajes y herramientas actuales y la variedad de modelos disponibles para ponerlos a prueba.

Es por ello que se identifican de forma explícita las siguientes limitaciones:

\begin{enumerate}
  \item \textbf{Expresividad limitada de los modelos de características disponibles:} Muchos de estos modelos se basan principalmente en relaciones estructurales básicas entre características. Aunque estos modelos permiten realizar análisis fundamentales, no siempre permiten representar escenarios más complejos donde intervienen elementos avanzados del lenguaje, lo que reduce la capacidad de explorar nuevas técnicas de razonamiento y análisis.
  \item \textbf{Complicaciones a la hora de evaluar y mejorar motores de razonamiento avanzados sobre modelos más complejos:} Esto es consecuencia directa del problema anteriormente descrito. Sin estos modelos más avanzados, la capacidad de validar nuevas aproximaciones se ve limitada, por lo que es más complicado estudiar hasta qué punto estos motores pueden aprovechar construcciones más avanzadas del lenguaje de modelado.
  \item \textbf{Necesidad de ampliar el conjunto de modelos disponibles para investigación en variabilidad software} La investigación en este ámbito requiere disponer de modelos que permitan experimentar con diferentes escenarios, niveles de complejidad y características del lenguaje. Sin una base amplia y diversa de modelos, resulta mucho más difícil comparar técnicas, evaluar resultados y avanzar en el desarrollo de herramientas más potentes.
\end{enumerate}

\section{Motivación}
\label{sec:motivacion}

\ph{Justifica por qué merece la pena resolver el problema y por qué con este enfoque.
Conecta el problema (sección anterior) con la solución que se va a proponer
(sin entrar todavía en cómo se implementa).}

\ph{Puntos típicos:
\begin{itemize}
  \item \ph{Por qué el enfoque elegido es adecuado.}
  \item \ph{Qué aporta respecto a soluciones existentes.}
  \item \ph{Encaje del trabajo en un contexto académico, profesional o de investigación
        (grupo de investigación, asignatura, empresa, publicación\dots).}
\end{itemize}}

La disponibilidad de modelos de características con distintos niveles de expresividad resulta fundamental para continuar avanzando en el desarrollo y evaluación de técnicas de análisis automático de variabilidad software. Como se ha comentado anteriormente, la evolución de los lenguajes de modelado permite representar escenarios cada vez más complejos, pero esta capacidad debe estar acompañada de modelos que permitan estudiar y validar el comportamiento de las herramientas encargadas de procesarlos.

Abordar esta necesidad dentro de un entorno de investigación permite trabajar sobre problemas que aparecen de forma directa durante el desarrollo de nuevas técnicas y herramientas. En este sentido, el contexto proporcionado por TASOVA Plus y DiversoLab ofrece un marco adecuado para contribuir al avance de soluciones relacionadas con la ingeniería de líneas de productos software, la arquitectura software y el análisis de variabilidad, lo que además permite seguir principios asociados al software libre y la reutilización del conocimiento científico.

Desde una perspectiva académica y profesional, este Trabajo Fin de Grado representa también una oportunidad para aplicar los conocimientos adquiridos durante el grado en un proyecto software real. Durante su desarrollo, he tenido la oportunidad de trabajar dentro de un ecosistema ya existente, enfrentándome a aspectos habituales en la ingeniería del software, como la comprensión de código desarrollado previamente, la integración de nuevas funcionalidades, la documentación o la validación mediante pruebas automatizadas.

Además, la participación en un proyecto vinculado a investigación me ha permitido adquirir experiencia en un contexto diferente al de los desarrollos software tradicionales. La necesidad de analizar información técnica, consultar documentación especializada y tomar decisiones fundamentadas ha favorecido el desarrollo de habilidades relacionadas con la búsqueda autónoma de conocimiento y la resolución de problemas abiertos.


\section{Punto de partida y contribución propia}
\label{sec:punto_partida}

Este Trabajo Fin de Grado se desarrolla sobre un conjunto de componentes y avances
previos que forman parte del ecosistema Flamapy-UVLHub. Por este motivo, resulta
necesario delimitar qué elementos constituyen la base sobre la que se construye el
trabajo y cuáles corresponden a las contribuciones realizadas durante el desarrollo
del presente proyecto.

En primer lugar, el diseño conceptual del generador de modelos de características
no surge desde cero en este trabajo. La propuesta de generación de modelos UVL con
diferentes niveles de expresividad está basada en el enfoque descrito por
Sundermann et al. \cite{sundermann2024generating}, donde se presenta un generador
capaz de producir modelos de características aprovechando las construcciones
avanzadas del lenguaje UVL. Este trabajo previo establece las bases conceptuales
sobre las que se desarrolla la implementación realizada en este TFG.

Además, el proyecto parte de la infraestructura existente del ecosistema UVLHub y
Flamapy. UVLHub ya proporcionaba una plataforma para la gestión y manipulación de
modelos UVL, así como un módulo \textit{generator} inicial sobre el que se realiza
la integración de la nueva funcionalidad.

La siguiente tabla resume la procedencia de los principales elementos relacionados
con el desarrollo:

\begin{table}[H]
\centering
\begin{tabular}{p{4cm}p{9cm}}
\toprule
\textbf{Elemento} & \textbf{Procedencia} \\
\midrule

Diseño conceptual del generador de modelos UVL con expresividad avanzada &
Basado en el trabajo de Sundermann et al. \cite{sundermann2024generating}, donde
se define el enfoque general de generación de modelos de características con
soporte para las construcciones avanzadas del lenguaje UVL. \\

Ecosistema Flamapy y plataforma UVLHub &
Infraestructura previa sobre la que se desarrolla el trabajo. En Flamapy participan diferentes grupos de investigación, entre ellos los pertenecientes a las Universidades de Sevilla, Málaga, Graz y la UNED. En cuanto a UVLHub, ha sido desarrollado principalmente por DiversoLab. \\

Generador de modelos de características como herramienta compatible con Flamapy &
Implementación desarrollada principalmente en este Trabajo Fin de Grado a partir del enfoque conceptual diseñado por Sundermann et al \cite{sundermann2024generating}. Incluye
la definición del modelo de configuración, la generación de estructuras
jerárquicas, atributos, restricciones y la integración con las operaciones de
Flamapy. \\

Integración del generador dentro del módulo \textit{generator} de UVLHub &
Desarrollo realizado durante este TFG, adaptando la aplicación existente para
utilizar el nuevo componente generador independiente. \\

Nueva interfaz del \textit{wizard} del generador &
Aportación realizada por David Romero dentro del módulo
\textit{generator} de UVLHub. Incluye la ampliación del \textit{wizard} de cinco a seis
pasos, el panel lateral de seguimiento de la configuración actual y el componente
\textit{slider} para facilitar la definición de distribuciones. \\

Refactorización de determinados elementos de UVLHub &
Trabajo realizado por David Romero dentro del proyecto global, incluyendo la
reorganización de algunos ficheros del módulo generador como
\textit{routes.py}, \textit{wizard.py} y otros elementos relacionados. Posteriormente, durante este TFG, se han realizado nuevas modificaciones arquitectónicas necesarias para integrar el generador.\\

Batería de pruebas unitarias del módulo generador &
Las pruebas relacionados con \textit{test\_engine.py},
\textit{test\_validators.py}, \textit{test\_wizard\_outputs.py} y
\textit{test\_wizard\_flow.py} fueron desarrollados en su mayoría por David Romero. Durante
el desarrollo del TFG algunas de estas pruebas han sido adaptados para ajustarse a
la nueva arquitectura del generador. \\

Generador basado en modelos de lenguaje (LLM) &
Desarrollado por David Romero dentro del ecosistema UVLHub. Aunque comparte
entorno con el presente trabajo, queda explícitamente fuera del alcance de este
Trabajo Fin de Grado. \\

\bottomrule
\end{tabular}
\caption{Procedencia de los principales elementos relacionados con el desarrollo.}
\label{tab:procedencia-elementos}
\end{table}


Por tanto, la contribución principal de este Trabajo Fin de Grado se centra en la
implementación del generador configurable de modelos de características como
componente independiente dentro de Flamapy y en su integración dentro de UVLHub.
En concreto, las principales aportaciones realizadas durante este trabajo son:

\begin{itemize}
    \item La implementación de la lógica de generación de modelos de características
    siguiendo el enfoque conceptual definido en trabajos previos, incorporando
    soporte para diferentes niveles de expresividad del lenguaje UVL.

    \item La creación de una arquitectura modular compatible con Flamapy,
    separando la configuración del generador, la representación del modelo y las
    operaciones necesarias para producir nuevos modelos.

    \item La implementación del sistema de configuración asociado al generador,
    permitiendo definir parámetros relacionados con características, jerarquía,
    atributos y restricciones.

    \item La integración del nuevo componente dentro del módulo \textit{generator} de UVLHub
    y la adaptación de la infraestructura existente para trabajar con esta nueva
    arquitectura.

    \item La validación del funcionamiento del sistema mediante pruebas
    automatizadas adaptadas a la implementación.
\end{itemize}

\section{Objetivos}
\label{sec:objetivos}

\ph{Define el objetivo general del trabajo en una o dos frases, y a continuación
los objetivos específicos en forma de lista. Formula cada objetivo específico
siguiendo los criterios \textbf{SMART}: \textbf{S}pecific (específico: una sola
cosa, sin ambigüedad), \textbf{M}easurable (medible: con una métrica o criterio
de éxito verificable), \textbf{A}chievable (alcanzable: realista con los recursos
y el plazo de un TFG), \textbf{R}elevant (relevante: conectado con una de las
limitaciones de la Sección~\ref{sec:problema}) y \textbf{T}ime-bound (acotado en
el tiempo: realizable dentro del calendario del Capítulo~\ref{chap:planificacion}).}

\ph{Usa verbos de resultado, no de actividad: evita ``estudiar'', ``aprender'' o
``profundizar en'' (no se puede verificar cuándo terminan) y prefiere
``implementar X que soporte Y'', ``alcanzar una cobertura del X\,\%'', ``reducir
el tiempo de Z por debajo de N segundos''. Ejemplo de objetivo bien formulado:
``Implementar un módulo de búsqueda que resuelva consultas sobre 10\,000
registros en menos de 200~ms, dando respuesta a la Limitación~1''.}

\ph{Objetivo general: ``Desarrollar/implementar/diseñar \dots con el fin de \dots''.}

\ph{\begin{enumerate}[label=\textbf{O\arabic*.}]
  \item \textbf{\ph{Objetivo específico 1:}} \ph{descripción concreta y medible.}
  \item \textbf{\ph{Objetivo específico 2:}} \ph{descripción concreta y medible.}
  \item \textbf{\ph{Objetivo específico 3:}} \ph{descripción concreta y medible.}
  \item \textbf{\ph{Objetivo específico 4:}} \ph{descripción concreta y medible.}
  \item \textbf{\ph{Objetivo específico 5 (opcional):}} \ph{descripción concreta y medible.}
\end{enumerate}
}

El objetivo principal de este trabajo es dar apoyo a la mejora de los motores de razonamiento existentes para el análisis de modelos de características mediante la obtención de una cantidad significativa y variada de modelos con expresividad avanzada. Como propuesta para conseguirlo, se pretende desarrollar una herramienta de generación automática de modelos de características aleatoria y completamente configurable, compatible con el ecosistema Flamapy, capaz de producir modelos UVL con diferentes niveles de expresividad \cite{sundermann2024generating}.

De forma más específica, podemos desgranar los objetivos de este trabajo de la siguiente forma:

\begin{enumerate}[label=\textbf{O\arabic*.}]
      \item Desarrollar un generador automático de modelos de características capaz de producir modelos UVL con diferentes niveles de complejidad y expresividad, ampliando la variedad de modelos disponibles para el análisis de variabilidad.

      \item Desarrollar un sistema de configuración flexible que permita modificar los parámetros principales del proceso de generación y adaptar los modelos obtenidos a distintos escenarios.

      \item Desarrollar una herramienta compatible con el ecosistema Flamapy e integrarla en UVLHub como solución para la generación configurable de modelos de características.

      \item Evaluar el comportamiento temporal del generador bajo escenarios de generación con un volumen elevado de modelos, analizando el coste asociado a la inicialización del entorno y al proceso de generación.

      \item Validar el funcionamiento del generador mediante una batería de pruebas automatizadas que permita comprobar la corrección de los modelos producidos y detectar posibles errores durante el proceso de generación.
\end{enumerate}

De forma complementaria, el trabajo persigue una serie de objetivos formativos
relacionados con la aplicación de conocimientos técnicos y metodológicos durante
el desarrollo. Su cumplimiento se evaluará de forma cualitativa a partir de las
evidencias documentadas en la memoria, incluyendo las decisiones de diseño, la
integración realizada, las pruebas desarrolladas y la planificación seguida.

\begin{enumerate}[resume,label=\textbf{O\arabic*.}]
      \item Reutilizar el conocimiento científico en el desarrollo de herramientas software (objetivo heredado de la red TASOVA Plus)
      \item Adquirir experiencia en el desarrollo y mantenimiento de proyectos software reales mediante la participación en un entorno de investigación y código abierto.
      \item Profundizar en el uso de tecnologías empleadas en el proyecto, especialmente Python, Flask y Docker, así como en la aplicación de pruebas automatizadas durante el desarrollo de un proyecto software integrado.
      \item Mejorar la capacidad de búsqueda, análisis y selección de información técnica relevante para resolver problemas dentro de un contexto de investigación.
      \item Aplicar buenas prácticas de planificación, gestión y documentación durante el desarrollo individual de un proyecto software de mayor duración y complejidad.
      
\end{enumerate}

\ph{Resume en la tabla siguiente el criterio de éxito de cada objetivo, cómo se
verificará y en qué capítulo. Esta tabla es la columna vertebral de la memoria:
el Capítulo~\ref{chap:resultados} dará el veredicto sobre \textbf{todos} los
criterios, recapitulando con \textbackslash ref la evidencia de aquellos que se
verifican en capítulos anteriores, y el Capítulo~\ref{chap:conclusiones}
sintetizará el grado de cumplimiento final. Si el criterio es cuantitativo,
fíjalo también como RNF en el Capítulo~\ref{chap:requisitos} con el
\textbf{mismo} umbral: un único umbral, dos vistas. Conviene rellenar la
columna «Cap.» con \textbackslash ref (p.\,ej.\ Capítulo~\ref{chap:pruebas})
para que la numeración se mantenga sola.}

\begin{longtable}{@{}P{1cm} P{6.2cm} P{4.6cm} P{1.7cm}@{}}
  \toprule
  \textbf{Obj.} & \textbf{Criterio de éxito (medible)} & \textbf{Cómo se verifica} & \textbf{Cap.} \\
  \midrule
  \endhead
  O1 & El generador permite producir modelos UVL automáticamente incorporando diferentes niveles de complejidad y expresividad del lenguaje, ampliando la variedad de modelos disponibles para el análisis de variabilidad. & Generación de modelos con configuraciones variadas y validación de que los modelos producidos cumplen la sintaxis y estructura esperada de UVL. & 6, 7 \\[0.3em]
  O2 & El generador dispone de un sistema de configuración que permite modificar los parámetros principales del proceso de generación y adaptar los modelos obtenidos a distintos escenarios. & Revisión de las opciones de configuración disponibles y ejecución de pruebas con diferentes configuraciones para comprobar su efecto sobre los modelos generados. & 6, 7 \\[0.3em]
  O3 & La solución desarrollada proporciona una herramienta compatible con el ecosistema Flamapy y puede ser utilizada desde UVLHub como herramienta de generación de modelos de características. & Verificación de la compatibilidad con los componentes empleados de Flamapy y comprobación del flujo completo de generación desde UVLHub. & 6, 7, 8 \\[0.3em]
  O4 & Evaluar el comportamiento temporal del generador en escenarios de generación con un volumen elevado de modelos de características. & Medición del tiempo de inicialización del generador y del tiempo requerido para generar un conjunto elevado de modelos UVL bajo una configuración representativa, analizando el coste asociado a cada fase del proceso. & 8 \\
  O5 & El generador cuenta con una batería de pruebas automatizadas que permita validar su correcto funcionamiento y detectar errores en los modelos producidos. & Ejecución de la suite de pruebas automatizadas y análisis de los resultados obtenidos. & 7 \\
  
  O6 & El desarrollo reutiliza conocimiento científico relacionado con la
  generación de modelos de características y lo aplica en una herramienta
  software concreta. & Análisis del trabajo de Sundermann et al. y trazabilidad
  de su enfoque conceptual en el diseño e implementación del generador. & 1, 2,
  5, 6 \\[0.3em]

  O7 & El desarrollo proporciona experiencia en la evolución de un proyecto
  software real, integrado en un ecosistema de investigación y código abierto.
  & Análisis de la integración del componente en Flamapy y UVLHub, del uso de
  Docker y de las pruebas ejecutadas sobre ambos proyectos. & 3, 5, 6, 7, 8
  \\[0.3em]

  O8 & El desarrollo aplica Python, Flask, Docker y pruebas automatizadas en la
  construcción e integración del generador. & Revisión de la implementación del
  componente, de su integración en UVLHub y de la ejecución de las suites de
  pruebas. & 5, 6, 7 \\[0.3em]

  O9 & Las decisiones técnicas del proyecto se apoyan en la búsqueda, análisis
  y selección de información técnica y científica relevante. & Revisión del
  estado del arte, de las fuentes consultadas y de las decisiones justificadas
  en el diseño e implementación. & 2, 5, 6 \\[0.3em]

  O10 & El desarrollo individual del proyecto aplica prácticas de planificación,
  gestión de riesgos y documentación técnica. & Revisión de la planificación,
  de la gestión de riesgos, de la documentación elaborada y de la reflexión
  final sobre el desarrollo. & 3, 9 \\[0.3em]
  \bottomrule
  \caption{Criterios de éxito de los objetivos y capítulos donde se verifican.}
  \label{tab:criterios-exito}
\end{longtable}

\section{Alcance}
\label{sec:alcance}

\ph{Delimita lo que el trabajo \textbf{sí} cubre y, sobre todo, lo que \textbf{no} cubre.
Es la sección que protege al autor de expectativas desmedidas durante la defensa.}

\ph{Una posible estructura:}

\ph{Lo que sí se aborda:}
\ph{
\begin{itemize}
  \item \ph{Funcionalidad / dominio / tipo de entradas que el sistema acepta.}
  \item \ph{Restricciones cuantitativas (rangos, tamaños, plataformas\dots).}
\end{itemize}
}

\ph{Quedan explícitamente fuera del alcance:}
\ph{
\begin{itemize}
  \item \ph{Funcionalidad excluida 1 (con una breve justificación).}
  \item \ph{Funcionalidad excluida 2 (con una breve justificación).}
  \item \ph{Funcionalidad excluida 3 (con una breve justificación).}
\end{itemize}
}

Este TFG comprende el desarrollo de un generador automático de modelos de características integrado dentro del ecosistema de herramientas asociado a Flamapy y orientado principalmente a su utilización desde UVLHub. La solución desarrollada permite generar modelos expresados en UVL mediante una configuración definida por el usuario, incorporando diferentes niveles de complejidad y expresividad con el objetivo de ampliar la disponibilidad de modelos para tareas de análisis e investigación en variabilidad software.

Dentro del alcance del proyecto se incluye el diseño e implementación del generador como herramienta compatible con Flamapy, siguiendo la arquitectura extensible del framework. Asimismo, se contempla su integración dentro del módulo \textit{generator} de UVLHub, permitiendo que los usuarios de la plataforma puedan interactuar con la funcionalidad desarrollada mediante una interfaz adaptada al flujo existente de la aplicación. Esta integración incluye la definición de los mecanismos necesarios para configurar el proceso de generación, la incorporación de soporte para diferentes construcciones del lenguaje UVL y la validación del funcionamiento del generador mediante pruebas automatizadas destinadas a comprobar la corrección de los modelos producidos.

En cuanto a lo que queda \textbf{fuera} del alcance del proyecto, se consideran los siguientes aspectos:

\begin{itemize}
    \item La creación de una aplicación web independiente para el generador. Aunque la integración realizada en UVLHub incorpora una interfaz que permite configurar y ejecutar el proceso de generación, el generador como componente reutilizable no incluye una interfaz propia. En caso de utilizarse fuera de UVLHub, sería necesario desarrollar una capa de interacción adaptada al contexto concreto de uso.
    
    \item La modificación o ampliación de otros módulos de UVLHub distintos al módulo \textit{generator}, ya que la intervención realizada se limita a la incorporación del generador dentro del módulo correspondiente.
    
    \item El desarrollo de nuevos motores de razonamiento o algoritmos de análisis de variabilidad. El objetivo del generador es proporcionar modelos con diferentes características y niveles de expresividad que puedan servir como base experimental para este tipo de herramientas, pero esto es algo que se pretende conseguir indirectamente con su uso.
\end{itemize}

\section{Estructura de la memoria}
\label{sec:estructura_memoria}

\ph{Cierra la introducción con un mapa de lectura: una o dos frases por capítulo.
No describas solo el contenido (``se presentan los requisitos''), sino el papel
de cada capítulo en el argumento de la memoria (``se establece qué debe hacer el
sistema, base contractual del diseño'').}

\ph{\begin{itemize}
  \item El \textbf{Capítulo~\ref{chap:estado_arte}} \ph{revisa los fundamentos
        teóricos y los trabajos relacionados, identificando el hueco que este
        trabajo cubre.}
  \item El \textbf{Capítulo~\ref{chap:planificacion}} \ph{describe la metodología,
        las fases, los riesgos y el presupuesto con los que se ha gestionado el
        proyecto.}
  \item El \textbf{Capítulo~\ref{chap:requisitos}} \ph{especifica qué debe hacer
        el sistema: requisitos, casos de uso y su trazabilidad, base contractual
        del resto del trabajo.}
  \item El \textbf{Capítulo~\ref{chap:diseno}} \ph{presenta cómo se estructura la
        solución: arquitectura, modelos y decisiones de diseño que responden a
        los requisitos.}
  \item El \textbf{Capítulo~\ref{chap:implementacion}} \ph{detalla cómo se
        materializa el diseño: tecnologías elegidas y módulos construidos.}
  \item El \textbf{Capítulo~\ref{chap:pruebas}} \ph{muestra cómo se ha verificado
        que el sistema cumple los requisitos y con qué estrategia de pruebas.}
  \item El \textbf{Capítulo~\ref{chap:resultados}} \ph{contrasta los resultados
        obtenidos con los criterios de éxito de la Tabla~\ref{tab:criterios-exito}
        y los discute frente a los trabajos relacionados.}
  \item El \textbf{Capítulo~\ref{chap:conclusiones}} \ph{da el veredicto sobre el
        cumplimiento de los objetivos, reconoce las limitaciones y propone líneas
        de trabajo futuro.}
\end{itemize}
}

% --- Transición al siguiente capítulo (último párrafo del roadmap) ---
\ph{Remata el roadmap con un párrafo puente de 2--4 líneas que lance al lector al
siguiente capítulo. Por ejemplo: «Definidos el problema, los objetivos y el
alcance, el Capítulo~\ref{chap:estado_arte} revisa los fundamentos y las
soluciones existentes sobre los que se construye la propuesta».}

El resto de la memoria se organiza como sigue:

\begin{itemize}
    \item El \textbf{Capítulo 2} presenta el \textbf{estado del arte} relacionado con la ingeniería de líneas de productos software, los modelos de características, el lenguaje UVL y las herramientas que forman parte del ecosistema en el que se desarrolla el trabajo. Además, se revisan trabajos relacionados con el objetivo de identificar las necesidades que justifican la propuesta desarrollada.

    \item El \textbf{Capítulo 3} describe la \textbf{planificación y gestión del proyecto}, incluyendo la metodología seguida durante el desarrollo, la organización de las tareas, la gestión de riesgos y los aspectos relacionados con la evolución del trabajo a lo largo del tiempo.

    \item El \textbf{Capítulo 4} define los \textbf{requisitos} que debe cumplir la solución desarrollada. En este capítulo se establecen las funcionalidades esperadas del generador, sus restricciones y la relación existente entre los requisitos y las pruebas destinadas a validarlos.

    \item El \textbf{Capítulo 5} presenta el \textbf{diseño de la solución} propuesta, describiendo la arquitectura general del sistema, la organización de sus componentes y las decisiones tomadas para integrar el generador dentro del ecosistema de Flamapy y UVLHub.

    \item El \textbf{Capítulo 6} detalla la \textbf{implementación} de la solución, incluyendo las tecnologías utilizadas, la estructura de la herramienta desarrollada, los mecanismos empleados para la generación de modelos y la integración realizada con la plataforma UVLHub.

    \item El \textbf{Capítulo 7} describe la estrategia de \textbf{pruebas} aplicada para verificar el correcto funcionamiento del sistema. En él se presentan las pruebas unitarias, de integración y de extremo a extremo utilizadas para validar la generación de modelos y la integración de la funcionalidad desarrollada.

    \item El \textbf{Capítulo 8} analiza los \textbf{resultados} obtenidos durante la evaluación del trabajo. Se contrastan los resultados con los criterios de éxito definidos para los objetivos y se discuten las aportaciones realizadas frente al contexto presentado en capítulos anteriores.

    \item El \textbf{Capítulo 9} recoge las \textbf{conclusiones} del trabajo, valorando el grado de cumplimiento de los objetivos planteados, las limitaciones encontradas y las posibles líneas de trabajo futuro derivadas de la solución desarrollada.
\end{itemize}

Definidos el contexto, el problema, los objetivos y el alcance del trabajo, el siguiente capítulo establece los fundamentos teóricos y tecnológicos necesarios para comprender el entorno de variabilidad software sobre el que se desarrolla la propuesta.

